\section{Coprime abelian canonical induction}
\label{sec:SF-coprime-abelian-canonical}

\paragraph{Status and provenance.}
Accepted proof for D-TD67 / Q2fw after complete stored-proof,
source-interface, code and first-output reading. Root and independent
whole-proof reviews and final documentary bindings pass. The exact-field canonical reduction
and normal-induction estimate are the accepted D-TD60 and D-TD56.
The parent asked whether those reductions constrain the remaining
proper-field source. The cyclic-prime, coprime-complement and actual
affine-factor argument below was derived by /root/sixfolds.
The original resumable note and the explicit unitary-qualification
precision are retained separately. No primitivity or quasiprimitivity
of the original coefficient representation is assumed.

\subsection*{The actual minimal source and statement}

Let \(K\) be imaginary quadratic. Let \(C\to\mathbf P^1\) be an actual
connected finite \(G\)-Galois cover in a family, with a faithful
absolutely irreducible \(K\)-representation \(\rho_K\) whose exact
complex character field is \(K\). Its multiplicity variation has
\(K\)-dimension six, signature \((3,3)\), and actual period rank nine.
Fix the signed determinant class of its naturally polarized primitive
factor, obtained from a Rosati-fixed rank-one idempotent, in
\(\mathbf Q^\times/N_{K/\mathbf Q}(K^\times)\).
Assume that the coefficient degree is minimal among all actual
sources with these exact field, signature, rank and natural determinant
data, including imprimitive sources. This is precisely the minimal
comparison class in D-TD56 and D-TD60.

Let \(A\) be an abelian minimal normal subgroup of \(G\) such that
\(\rho_{\mathbf C}|_A\) is not homogeneous. Choose a canonical
\(A\)-isotypic block \(B\), put \(H=\operatorname{Stab}_G(B)\),
and let \(\tau\) be the \(H\)-representation on \(B\). Assume
\(e=\dim_{\mathbf C}B\ge2\).

\begin{theorem}[D-TD67: coprime abelian canonical branch]
The subgroup \(A\) is cyclic of odd prime order \(p\). With
\(L=\mathbf Q(\chi_\tau)\) and \(k=[G:H]>1\), one has
\[
 A=C_p,\qquad H=C_G(A)\triangleleft G,\qquad
 L=K(\zeta_p),\qquad
 G/H\simeq\operatorname{Gal}(K(\zeta_p)/K),\qquad
 \ker\tau=1.
 \tag{CAC.1}
\]
These assertions do not require either coprimality condition below.
If, in addition,
\[
             p\nmid |H/A|,\qquad p\nmid\operatorname{disc}(K),
 \tag{CAC.2}
\]
then
\[
                    p\le83,\qquad k=p-1\le82.
 \tag{CAC.3}
\]
In this coprime branch there is a characteristic complement \(D\) in
\(H\), normal in \(G\), and an actual complex tensor factorization of
the original coefficient representation
\[
 H=A\times D,\qquad
 \rho_{\mathbf C}=W_\eta\otimes U,\qquad
 \dim W_\eta=e,\quad \dim U=k,\qquad
 P_U\simeq G/D\simeq C_p\rtimes\mathbf F_p^\times.
 \tag{CAC.4}
\]
Here \(P_U\) is the actual factorwise projective image of \(G\) on
\(U\). No model of \(\eta\) over \(K\), and no Schur-index-one
assertion for \(\eta\), is needed.
\end{theorem}

\subsection*{1. The faithful prime character and exact field}

An abelian finite minimal normal subgroup is elementary abelian:
its primary subgroups are characteristic, so it is a \(p\)-group,
and the characteristic subgroup \(pA\), being proper, is zero.
Write \(A=(\mathbf F_p)^d\). Its irreducible complex characters
are linear. Let \(\theta\) label \(B\). It is nontrivial: if it
were trivial, all its conjugates would be trivial, contrary to
faithfulness of \(\rho\).

Apply the complete canonical normal-isotypic corollary D-TD60,
Section 4 and (GSB.12), including its minimal-source conclusion.
It gives
\[
 \rho_{\mathbf C}=\operatorname{Ind}_H^G\tau,\qquad
 H\triangleleft G,\qquad
 L=K\mathbf Q(\theta),\qquad
 k=[L:K],\qquad
 \ker\tau=1.
 \tag{CAC.5}
\]
The \(k\) distinct conjugates of \(\tau\) are its full regular
\(\operatorname{Gal}(L/K)\)-orbit. Restricting to \(A\) shows that
all canonical block characters are Galois conjugates of \(\theta\).
An automorphism of a cyclotomic field preserves the kernel of a
linear character. Thus these characters have the same kernel.
Their common kernel is \(\ker(\rho|_A)=1\). Consequently \(\theta\)
is faithful and embeds \(A\) into \(\mathbf C^\times\).
An elementary abelian subgroup of \(\mathbf C^\times\) is cyclic,
so \(d=1\) and \(A=C_p\).

The stabilizer of the faithful character \(\theta\), and therefore
of its canonical isotypic block, is exactly \(C_G(A)\).
Its character field is \(\mathbf Q(\zeta_p)\), proving (CAC.1).
If \(p=2\), then \(L=K\) and (CAC.5) gives \(k=1\), contrary
to nonhomogeneity. Hence \(p\) is odd.
Notice that the use of equal kernels comes after the complete
Galois-orbit conclusion; it is not inferred merely from normality
of \(H\).

Put \(f=|\operatorname{disc}(K)|\). The exact quadratic
cyclotomic inclusion is \(K\subseteq\mathbf Q(\zeta_f)\);
see Milne, \cite{HD:MilneCFT2020}, V Example 3.11, printed p.~159.
The cyclotomic Galois identification in V Example 3.2, p.~157,
together with the Chinese remainder theorem, gives
\[
 \operatorname{Gal}(\mathbf Q(\zeta_{pq})/\mathbf Q)
 \simeq
 \operatorname{Gal}(\mathbf Q(\zeta_p)/\mathbf Q)
 \times\operatorname{Gal}(\mathbf Q(\zeta_q)/\mathbf Q)
 \quad\text{if }\gcd(p,q)=1.
 \tag{CAC.6}
\]
This also follows directly from
\(\varphi(pq)=\varphi(p)\varphi(q)\) and
\(\mathbf Q(\zeta_{pq})=\mathbf Q(\zeta_p,\zeta_q)\).
Apply it with \(q=f\), using (CAC.2). Then \(k=p-1\);
the conjugation action \(G/H\) on \(A\) is the full
\(\mathbf F_p^\times\).

The accepted cyclic estimate D-TD56, (NID.4), now applies to the
same original source. In particular, with the original genuine
inertia matrices \(T_i\), it supplies
\[
 e\ge5k-3,\qquad
 \dim\rho=ek,\qquad
 \sum_i\operatorname{rank}(T_i-I)=2ek+6.
 \tag{CAC.7}
\]
Both quotient bases occurring in that normal-induction theorem
are rational. Its mark and rank argument is an actual-source
statement. We have not replaced the source, enlarged \(K\), or
changed the natural determinant in deriving (CAC.7).

\subsection*{2. A unique central complement and its invariant type}

The subgroup \(A=C_p\) is central in \(H\) and is its unique
Sylow \(p\)-subgroup under (CAC.2). We prove the required splitting
and uniqueness explicitly. Write \(A\) additively and put \(Q=H/A\).
Choose a normalized section \(s:Q\to H\), with central cocycle
\(c:Q\times Q\to A\), so that
\(s(q)s(r)=c(q,r)s(qr)\). Define
\[
 a(q)=|Q|^{-1}\sum_{z\in Q}c(q,z)\in A,\qquad
 c(q,r)=a(q)+a(r)-a(qr).
 \tag{CAC.8}
\]
The inverse exists because \(p\nmid|Q|\). The second equality
follows by summing
\(c(q,r)+c(qr,z)=c(r,z)+c(q,rz)\) over \(z\).
Multiplying \(s(q)\) by the central element \(-a(q)\) kills the
cocycle. The adjusted section is a homomorphism and its image is
a complement \(D\), so \(H=A\times D\).

Every other complement is the graph of a homomorphism \(D\to A\);
that homomorphism is trivial because their orders are coprime.
Thus \(D\) is unique. As \(A\) is the unique Sylow \(p\)-subgroup
of \(H\), both \(A\) and the unique complement \(D\) are
characteristic in \(H\). In particular \(D\triangleleft G\).

Since \(A\) acts through \(\theta\) on \(B\), one can write
\(\tau=\theta\otimes\eta\) on \(A\times D\), where \(\eta\) is
an irreducible \(D\)-representation of degree \(e\).
It is faithful by \(\ker\tau=1\).
Let \(t=\exp D\), so \(\gcd(p,t)=1\). Character values of \(\eta\)
are sums of \(t\)-th roots of unity and also values of \(\chi_\tau\)
on \(D\). Hence
\[
       \mathbf Q(\chi_\eta)\subseteq
       L\cap\mathbf Q(\zeta_t),\qquad
       \mathbf Q(\chi_\eta)\subseteq K.
 \tag{CAC.9}
\]
For the second inclusion, take any
\(\sigma\in\operatorname{Gal}(L/K)\).
By (CAC.6) with \(q=ft\), its action on the \(p\)-th roots
extends to an automorphism of \(\mathbf Q(\zeta_{pft})\) fixing
\(\mathbf Q(\zeta_{ft})\). This extension agrees with \(\sigma\)
on \(L=K(\zeta_p)\), and fixes every value of \(\chi_\eta\).
The fixed field in \(L\) is \(K\), proving (CAC.9).

For every \(g\in G\), (CAC.5) identifies \(\tau^g\) with some
\(\tau^\sigma\). Restrict to the actual normal subgroup \(D\).
Equation (CAC.9) gives
\(\chi_{\eta^g}=\chi_{\eta^\sigma}=\chi_\eta\), so
\(\eta^g\simeq\eta\). This is invariance of a complex irreducible
type. It does not assert descent of that type to a \(K\)-model.

\subsection*{3. The factorwise quotient is the full affine group}

Mackey restriction to the normal subgroup \(H\), followed by
restriction to \(D\), gives
\(\rho|_D\simeq\eta^{\oplus k}\).
Evaluation therefore realizes the original space as
\(V=W_\eta\otimes U\), where
\(U=\operatorname{Hom}_D(W_\eta,V)\) has dimension \(k\).
The algebra generated by \(D\) is
\(\operatorname{End}(W_\eta)\otimes I\); it is normalized by \(G\).
An automorphism of a full matrix algebra is inner, as is seen
by transporting its standard irreducible module and matrix units.
Thus every \(g\in G\) has a factorization \(A_g\otimes B_g\),
unique up to reciprocal scalars. Multiplication consequently
defines actual homomorphisms from \(G\) to the two projective
linear groups. The second has finite image \(P_U\), and kills \(D\).

Put \(F=G/D\). The images of \(A\) and \(H\) give an exact sequence
with kernel \(A=C_p\) and quotient \(G/H=C_k\).
Moreover \(C_F(A)=A\): a commutator of a lift with \(A\)
lies in \(A\), so being trivial modulo \(D\) makes it trivial,
since \(A\cap D=1\). Lift a generator of \(C_k\) to \(x\in F\).
Then \(x^k\in A\) is fixed by \(x\). The corresponding nontrivial
multiplier in \(\mathbf F_p^\times\) fixes no nonidentity element
of \(A\), so \(x^k=1\). Therefore
\[
                 F=C_p\rtimes C_{p-1}.
 \tag{CAC.10}
\]

The restriction of \(A\) to \(U\) has all the nontrivial
characters of \(C_p\), each with multiplicity one.
Indeed these are the original canonical block characters, and
each original block has the common \(D\)-degree \(e\).
Since \(p>2\), no nonidentity element of \(A\) acts scalarly on \(U\).
The kernel of \(F\to P_U\) is normal and intersects \(A\) trivially.
It centralizes \(A\), because its commutators with \(A\) lie in
that intersection. It is consequently contained in \(C_F(A)=A\)
and is trivial. This proves the actual quotient assertion in
(CAC.4); no quotient of an auxiliary inducing group has been
substituted for a quotient of the original \(G\).

Choose invariant positive Hermitian forms. The \(D\)-invariant
form on \(V=W_\eta\otimes U\) is a tensor product of a form on
the irreducible \(W_\eta\) and a form on its multiplicity space.
Since \(G\) acts unitarily, the factors \(A_g,B_g\) can be made
unitary by reciprocal scalar rescaling. For a nontrivial
translation \(a\in C_p\), every unitary representative of its
projective action has trace, up to a unit scalar,
\(\sum_{j=1}^{p-1}\zeta_p^{aj}=-1\).
For an element outside \(A\), conjugation permutes these
character lines by a multiplier different from one. There is
no fixed line, so the trace is zero. Thus
\[
 \frac{|\operatorname{tr}B|}{k}\le\frac1k
 \quad\text{for every unitary representative \(B\) of a
 nonidentity element of \(P_U\)}.
 \tag{CAC.11}
\]
This qualifier is essential: no absolute trace assertion is
made for an arbitrarily rescaled nonunitary representative.

\subsection*{4. The actual quotient curve is hyperbolic for large \(p\)}

For contradiction suppose \(p>83\). Let
\(Z=C/\ker(G\to P_U)\). This is an actual connected
\(F\)-Galois cover of the original sphere.
If \(Z\) had genus zero, \(F\) would act faithfully in
\(\operatorname{PGL}_2(\mathbf C)\).
A nontrivial finite-order element of order \(p\) is conjugate
to \(z\mapsto\zeta_p z\) and has exactly two fixed points.
Its normalizer permutes those two points. The subgroup fixing
both individually consists of scalings and centralizes it.
Thus the conjugation image of the normalizer on \(C_p\)
has order at most two. This contradicts the full image
\(\mathbf F_p^\times\) of order \(p-1>2\).

If \(Z\) had genus one, choose an origin and regard it as an
elliptic curve \(E\). The complete statements and proofs of
\cite{TD:MilneEC2021}, IV Proposition 7.13 and Remark 7.18,
printed pp.~140--141, give
\[
 \operatorname{Aut}(E)
   =E(\mathbf C)\rtimes\operatorname{Aut}(E,0),
 \qquad
 |\operatorname{Aut}(E,0)|\in\{2,4,6\}.
 \tag{CAC.12}
\]
Since \(p>3\), \(A=C_p\) lies in the translation subgroup.
Conjugation of translations by translations is trivial;
its conjugation image under \(F\) factors through
\(\operatorname{Aut}(E,0)\). That image has order at most six,
again contradicting the faithful image of order \(p-1\).
Thus \(g(Z)\ge2\).

Let \(o_i\) be the orders of the images in \(P_U\) of the original
inertia elements, retaining \(o_i=1\) at omitted quotient branch
points. Riemann--Hurwitz on this actual quotient says
\[
 \sigma=\sum_i(1-1/o_i)>2,\qquad
                    \sigma\ge2+\frac1{42}.
 \tag{CAC.13}
\]
Here is the elementary gap argument, to specify every order case.
Discard the entries equal to one and let their number be \(s\).
For \(s\ge5\) the excess over two is at least \(1/2\).
For \(s=4\), all four orders two give excess zero; otherwise
the reciprocal sum is at most \(3/2+1/3\), giving excess at least
\(1/6\). For \(s=3\), sort the orders \(a\le b\le c\).
If \(a\ge4\) the excess is at least \(1/4\).
If \(a=3,b=3\), positivity forces \(c\ge4\) and gives \(1/12\);
if \(a=3,b\ge4\), it gives at least \(1/6\).
If \(a=2,b=2\), positivity is impossible.
If \(a=2,b=3\), positivity forces \(c\ge7\), giving \(1/42\).
If \(a=2,b=4\), it forces \(c\ge5\), giving \(1/20\);
if \(a=2,b\ge5\), the excess is at least \(1/10\).
Finally \(s\le2\) cannot have positive excess.
This proves (CAC.13) without assuming a presentation of \(F\).

\subsection*{5. Full original periods and the source budget}

Fix one original inertia \(T_i\), of its full coefficient order
\(L_i\). Choose unitary factors \(T_i=A_i\otimes B_i\).
The equality \(T_i^{L_i}=I\) implies
\(A_i^{L_i}=\alpha_i I\) and
\(B_i^{L_i}=\alpha_i^{-1}I\), with \(|\alpha_i|=1\).
Choose a unit scalar \(\mu_i\) with
\(\mu_i^{L_i}=\alpha_i^{-1}\).
Replacing the factors by \(\mu_i A_i,\mu_i^{-1}B_i\)
makes both \(L_i\)-th powers equal to the identity.
Their product remains the original \(T_i\).

The projective order \(o_i\) divides \(L_i\).
Exactly \(L_i/o_i\) exponents in a full period have scalar
\(B_i^j\), with normalized absolute trace one.
All other exponents satisfy (CAC.11).
The entire other factor has normalized absolute trace at most
one, so Fourier inversion for every target eigenvalue gives
\[
 \frac{\operatorname{rank}(T_i-\lambda I)}{ek}
 \ge 1-\frac1{o_i}-\frac{1-1/o_i}{k}
 =(1-1/o_i)(1-1/k).
 \tag{CAC.14}
\]
Explicitly, its target multiplicity is
\(L_i^{-1}\sum_{j=0}^{L_i-1}\lambda^{-j}
\operatorname{tr}(T_i^j)\), and its absolute value is at most
the sum of the absolute traces just bounded.
Targets outside the \(L_i\)-th roots have multiplicity zero.
This is the full original period, not just one chosen period
in the quotient. Scalar quotient entries and the complete
partner \(W_\eta\) have both been retained.

Apply (CAC.14) at \(\lambda=1\), sum over the original genuine
inertia, and use (CAC.7) and (CAC.13). One obtains
\[
 2+\frac6{ek}
 \ge(1-1/k)\left(2+\frac1{42}\right),
 \qquad\text{hence}\qquad e(k-85)\le252.
 \tag{CAC.15}
\]
For \(k\ge86\), (CAC.7) instead gives
\[
 e(k-85)\ge(5k-3)(k-85)\ge427>252.
 \tag{CAC.16}
\]
The polynomial is increasing on that integer range.
This contradiction under \(p>83\) proves \(k\le85\) whenever
that case could occur. Since \(k=p-1\) and \(p\) is an odd prime,
the primes above \(83\) have \(k\ge88\).
Thus every source in the coprime branch has \(p\le83\),
and then \(k=p-1\le82\), proving (CAC.3).

\paragraph{Exact boundary of the result.}
The cyclic-prime and exact-field assertions (CAC.1) hold for the
canonical nonhomogeneous abelian minimal-normal system without
(CAC.2). The numerical conclusion uses both coprimality assumptions.
The cases \(p\mid|H/A|\), ramified \(p\), and nonabelian minimal
normal subgroups have not been excluded by this proof.
For a fixed \(K\), its ramified primes form a separate finite list;
no stronger ramified-prime conclusion is needed here.
No representation \(\eta\) has been assumed realizable over \(K\);
there is no enlargement of the original coefficient field or change
of its natural determinant. The case \(e=1\) is outside the
normal-induction rational-base argument and is not asserted here.
The separate accepted monomial-source result applies only with its
own exact naturally nonsplit hypotheses.
This is a restriction on actual minimal sources, not a construction
of such sources and not an automatic reduction of every minimal
source to a quasiprimitive one.
